\section*{Bab \ref{ch:g9:landforms} --- Bentuk Lahan}

\begin{solution}{exo:g9:landforms:1}
(a) Angin. (b) Air yang melarutkan \omterm{def:g3:stones-pebbles-sand:rock}{batuan} (dan mengendapkan kalsit di dalam
\omterm{def:g9:landforms:karst}{gua}). (c) \omterm{def:g2:water-on-earth:ocean}{Laut} (\omterm{def:g9:landforms:beach}{hanyutan sepanjang pantai}). (d) Es. (e) \omterm{def:g2:water-on-earth:ocean}{Laut}. (f) \omterm{def:g1:landscapes-around-us:water}{Sungai}.
\end{solution}

\begin{solution}{exo:g9:landforms:2}
(a) \omterm{def:g9:landforms:uplift}{Pengangkatan}. (b) \omterm{def:g9:landforms:dune}{Lereng luncur}. (c) \omterm{def:g9:landforms:karst}{Dolina}. (d) \omterm{def:g9:landforms:beach}{Hanyutan sepanjang
pantai}. (e) \omterm{def:g9:landforms:karst}{Karst}.
\end{solution}

\begin{solution}{exo:g9:landforms:3}
(a) Salah: tanduk-tanduknya menunjuk searah angin; hidungnya yang bundar
menghadap angin. (b) Benar. (c) Salah: angin memindahkannya butir demi
butir, menaiki lereng yang landai dan menuruni \omterm{def:g9:landforms:dune}{lereng luncur}. (d) Benar.
\end{solution}

\begin{solution}{exo:g9:landforms:4}
Urutannya: (b), (d), (e), (a), (c).
\end{solution}

\begin{solution}{exo:g9:landforms:5}
(a) Dari barat: tanduk-tanduknya menunjuk searah angin. (b) Di sisi timur,
di antara kedua tanduk. (c) $60 \div 6 = 10$ tahun. (d) \omterm{def:g9:landforms:dune}{Gumuk} yang lebih
besar mengandung lebih banyak \omterm{def:g3:stones-pebbles-sand:sand}{pasir}, dan semua \omterm{def:g3:stones-pebbles-sand:sand}{pasir} itu harus diangkut
melewati puncaknya agar \omterm{def:g9:landforms:dune}{gumuk} itu dapat bergerak maju.
\end{solution}

\begin{solution}{exo:g9:landforms:6}
(a) $7 \times 1 = \qty{7}{m}$. (b) $2 \div 12 = \frac{1}{6}$ tahun.
(c) Mundurnya tujuh tahun dalam $\frac{1}{6}$ tahun: $7 \times 6 = 42$
kali lebih cepat. (d) Ombak mula-mula memahat takik; \omterm{def:g3:stones-pebbles-sand:rock}{batuan} di atasnya
menggantung di atas takik itu untuk sementara waktu, tanpa bergerak; lalu
bagian yang menggantung itu runtuh sekaligus, dan tepi \omterm{def:g9:landforms:cliff}{tebing} melompat
mundur.
\end{solution}

\begin{solution}{exo:g9:landforms:7}
(a) Ke arah timur. (b) Tiga gerakan naik ke \omterm{def:g9:landforms:beach}{gisik} yang menyerong ke arah
timur laut, masing-masing diikuti gerakan lurus kembali turun ke arah
selatan: zigzag yang bergerak ke timur. (c) \omterm{def:g9:landforms:beach}{Lidah pasir}, yang tumbuh ke
arah timur melintangi mulut teluk. (d) \omterm{def:g9:landforms:beach}{Lidah pasir} melindungi air di
belakangnya dari ombak; di air yang tenang, \omterm{def:g3:stones-pebbles-sand:sand}{lumpur} yang halus dapat
mengendap.
\end{solution}

\begin{solution}{exo:g9:landforms:8}
(a) Air \omterm{def:g4:clouds-rain-snow:precipitation}{hujan} menyerap karbon dioksida dari udara dan dari tanah. (b) Asam
lemah melarutkan batu kapur secara perlahan, tetapi hampir tidak
menyerang granit. (c) Air hanya dapat masuk ke dalam \omterm{def:g3:stones-pebbles-sand:rock}{batuan} melalui
retakannya dan di sepanjang lapisannya; air melebarkan celah-celah ini
lebih dulu. (d) \omterm{def:g3:stones-pebbles-sand:rock}{Batuan} di bawah permukaan larut dan tanah amblas
perlahan-lahan; atau atap sebuah \omterm{def:g9:landforms:karst}{gua} runtuh.
\end{solution}

\begin{solution}{exo:g9:landforms:9}
(a) Ketika tetes air bertemu dengan udara \omterm{def:g9:landforms:karst}{gua}, tetes itu melepaskan
sebagian karbon dioksidanya dan tidak dapat lagi menahan semua batu kapur
yang terlarut, yang keluar sebagai kalsit. (b) Stalagmit tumbuh dari tetes
yang memercik di lantai lalu menyebar; tidak ada tetes yang menggantung di
ujungnya untuk membangun tabung. (c) $\qty{30}{cm} = \qty{300}{mm}$;
$300 \div 0.1 = 3000 = 3 \times 10^3$ tahun.
\end{solution}

\begin{solution}{exo:g9:landforms:10}
(a) $4600 \div 10^7 = 4.6 \times 10^{-4}$~m setahun, atau \qty{0.46}{mm}
setahun. (b) $0.46 - 0.2 = \qty{0.26}{mm}$ setahun. (c) $\qty{1}{m} =
\qty{1000}{mm}$; $1000 \div 0.46 \approx 2200$, sekitar dua ribu kali
lebih cepat. (d) \omterm{def:g9:landforms:uplift}{Pengangkatan} itu berlangsung terus selama jutaan tahun:
$\qty{0.26}{mm} \times
10^7 = 2.6 \times 10^6$~mm, sekitar \qty{2.6}{km}.
\end{solution}

\begin{solution}{exo:g9:landforms:11}
(a) Bentuk: \omterm{def:g9:landforms:dune}{gumuk pasir} bulan sabit dengan \omterm{def:g9:landforms:dune}{lereng luncur} di sisi tenggara.
Tenaga dan proses: angin, yang bertiup dari barat laut, \omterm{def:g5:water-shapes-the-land:deposition}{pengendapan}.
Bahan: \omterm{def:g3:stones-pebbles-sand:sand}{pasir} lepas. (b) Bentuk: \omterm{def:g9:landforms:cliff}{tebing} di atas paparan datar yang
tersingkap saat air surut. Tenaga dan proses: \omterm{def:g2:water-on-earth:ocean}{laut}, \omterm{def:g5:water-shapes-the-land:erosion}{erosi}; paparan itu
adalah dataran abrasi yang ditinggalkan \omterm{def:g9:landforms:cliff}{tebing} yang mundur. Bahan: \omterm{def:g3:stones-pebbles-sand:rock}{batuan}
putih, kemungkinan kapur, \omterm{def:g3:stones-pebbles-sand:rock}{batuan} yang lunak. (c) Bentuk: \omterm{def:g1:landscapes-around-us:water}{sungai} kecil yang
menghilang ke dalam lubang, \omterm{def:g6:water-underground:spring}{mata air} yang jauh. Tenaga dan proses: air
yang melarutkan \omterm{def:g3:stones-pebbles-sand:rock}{batuan} dan mengalir di bawah tanah. Bahan: batu kapur;
dataran tinggi itu adalah \omterm{def:g9:landforms:karst}{karst}.
\end{solution}

\begin{solution}{exo:g9:landforms:12}
(a) \omterm{def:g9:landforms:dune}{Gumuk} yang kecil: $45 \div 6 = 7.5$ tahun; \omterm{def:g9:landforms:dune}{gumuk} yang lebih besar: $45 \div 3 =
15$ tahun. (b) $45 \div 100 = 0.45$ tahun, sekitar lima atau enam bulan.
(c) Misalnya, menanam rumput di \omterm{def:g9:landforms:dune}{gumuk-gumuk} itu: akar-akarnya menahan
\omterm{def:g3:stones-pebbles-sand:sand}{pasir}, sehingga angin tidak dapat lagi menerbangkannya.
\end{solution}

\begin{solution}{pb:g9:landforms:1}
\textbf{1.} $46.0 - 30.4 = \qty{15.6}{m}$.

\textbf{2.} $15.6 \div 20 = \qty{0.78}{m}$ setahun.

\textbf{3.} Antara 2012 dan 2014: $41.0 - 36.2 = \qty{4.8}{m}$ dalam dua
tahun, atau $4.8 \div 2 = \qty{2.4}{m}$ setahun.

\textbf{4.} $2.4 \div 0.78 \approx 3$ kali lebih cepat. Kemungkinan ada
musim dingin yang penuh \omterm{def:g5:natural-hazards:storm}{badai}, ketika bagian-bagian besar yang menggantung
runtuh.

\textbf{5.} \qty{78}{cm} setahun; $0.78 \times 100 = \qty{78}{m}$ per
abad.

\textbf{6.} Ya: \qty{0.78}{m} setahun sedikit kurang dari satu meter setahun.

\textbf{7.} $17 \div 0.78 \approx 22$ tahun.

\textbf{8.} \omterm{def:g9:landforms:cliff}{Tebing} mundur dengan lompatan: selang waktu dua tahun mungkin
memuat satu runtuhan besar, atau tidak ada sama sekali. Selama 20 tahun,
lompatan-lompatan dan tahun-tahun yang tenang saling merata.

\textbf{9.} Tepi \omterm{def:g9:landforms:cliff}{tebing} masih harus menempuh $30.4 - 10 = \qty{20.4}{m}$:
$20.4 \div 0.78
\approx 26$ tahun, jadi pada tahun 2050.

\textbf{10.} $2 \times 0.78 = \qty{1.56}{m}$ setahun; $30.4 \div 1.56
\approx 19.5$ tahun setelah 2024, jadi pada tahun 2043.

\textbf{11.} Tepi \omterm{def:g9:landforms:cliff}{tebing} dapat runtuh kapan saja; berdiri di atasnya, atau
di dekatnya, berbahaya.

\textbf{12.} $30.4 \div 0.78 \approx 39$ tahun setelah 2024: pada tahun
\textbf{2063}.
\end{solution}
